% !TEX TS-program = xelatex
% !TEX encoding = UTF-8 Unicode
% !Mode:: "TeX:UTF-8"

\documentclass{resume}
\usepackage{zh_CN-Adobefonts_external} % Simplified Chinese Support using external fonts (./fonts/zh_CN-Adobe/)
%\usepackage{zh_CN-Adobefonts_internal} % Simplified Chinese Support using system fonts
\usepackage{linespacing_fix} % disable extra space before next section
\usepackage{cite}
\usepackage{geometry}
\usepackage{multicol}
\geometry{top=0.40in,bottom=0.40in}
\setlist[itemize]{topsep=0.10em,parsep=0.1ex,leftmargin=1.5pc}
\linespread{0.85}

\begin{document}
\pagenumbering{gobble} % suppress displaying page number

\name{王海翔}

% {E-mail}{mobilephone}{homepage}
% be careful of _ in emaill address
\contactInfo{(+86) 156-6207-7668}{jt846962421@gmail.com}{AI Agent研发工程师}{GitHub: @JackWPP}
% {E-mail}{mobilephone}
% keep the last empty braces!
%\contactInfo{xxx@yuanbin.me}{(+86) 131-221-87xxx}{}
 
\section{个人总结}
本人具备扎实的计算机专业基础与AI工程实践能力，专注于大模型应用与多Agent系统研发。在校期间主导多个创新创业项目，擅长RAG技术、提示词工程及私有化部署。曾带领团队获"挑战杯"北京市一等奖等多项荣誉,\textbf{具备优秀的团队协作、危机处理与项目管理能力，现已保研至本校化工学院攻读硕士学位。}

% \section{\faGraduationCap\ 教育背景}
\section{教育背景}
\datedsubsection{\textbf{北京化工大学},计算机科学与技术,\textit{工学学士}}{2021.9 - 2026.6}
\ 获\textbf{挑战杯北京市一等奖}、\textbf{计算机能力挑战赛国家级二等奖}、\textbf{iCAN大赛北京市二等奖}等荣誉，预计2026年6月毕业
\datedsubsection{\textbf{北京化工大学},化学工程与技术,\textit{工学硕士}}{2026.9 - 2029.6}
\ 已保研至本校化学学院\textbf{化学工程与技术}专业，研究方向：\textbf{基于计算机视觉的电解水电极表面气泡追踪与量化分析}


% \section{\faCogs\ IT 技能}
\section{技术能力}
% increase linespacing [parsep=0.5ex]
\begin{itemize}[parsep=0.2ex]
  \item \textbf{编程语言}: Python, C/C++, Java, Shell, SQL
  \item \textbf{AI/深度学习}: PyTorch, RAG技术, Prompt Engineering, Dify框架, 多Agent系统
  \item \textbf{开发框架}: FastAPI, React, Flask
  \item \textbf{数据库与工具}: MySQL, MongoDB, Git, Linux
  \item \textbf{专业技能}: 计算机视觉, 大模型应用开发, 私有化部署, 多模态数据处理
\end{itemize}

% \end{itemize}

\section{校园实践}
\datedsubsection{\textbf{北京化工大学 | 北校区工作办公室}, 特邀摄影师}{2023 - 2025}
\begin{itemize}[parsep=0.1ex]
  \item 负责学校重要活动现场拍摄与视觉记录，包括重要来宾跟拍、活动全程纪实等，确保影像素材的专业性与完整性
  \item 连续两年负责新生入学"第一张照片"拍摄及宣传片制作，作品获校官方媒体刊发，单条播放量突破\textbf{10万+}
\end{itemize}


\section{项目经历}
\datedsubsection{\textbf{FoxSay 狐说 | 课程级AI学习Copilot},全栈开发(独立)}{2025 - 至今}
\begin{itemize}[parsep=0.05ex]
  \item \textbf{CRAG边界控制：}自研Confidence-RAG三档门控——\texttt{$\ge$\,0.72}附页码引用、\texttt{0.55--0.72}标注基于通用知识、\texttt{<\,0.55}直接拒答，解决通用AI\textbf{幻觉与超纲}问题
  \item \textbf{全栈实现：}前端 \textbf{React 19 + TS 6 + Tailwind v4 + Vite 8}（\texttt{reactflow}+\texttt{@dagrejs/dagre} 知识图谱、\texttt{katex} 公式）；后端 Python + \textbf{LangGraph} 多Agent，\textbf{168 个后端测试用例全部通过}
  \item \textbf{完整闭环：}4阶段Pipeline构建知识体系 $\to$ 7工具ReAct混合检索 $\to$ React Flow知识图谱 $\to$ 9步状态机超级备考 $\to$ 多题型自动出题
  \item \textbf{狐狸人设：}文案"贱萌但靠谱"风格 + \texttt{fox-breathe}/\texttt{fox-shimmer}/\texttt{fox-stagger-in}/\texttt{fox-pulse-ring} 等命名CSS动画，视觉/文案/交互统一
\end{itemize}
\datedsubsection{\textbf{PCBTool.AI | 智能PCB设计辅助系统},算法负责人}{2025 - 至今}
\begin{itemize}[parsep=0.1ex]
  \item \textbf{技术架构：}基于大模型+RAG构建多模态智能系统，支持自然语言/电路图多类型输入，自动生成项目文档、部署指南及PCB Layout方案
  \item \textbf{部署优化：}主导Dify私有化部署，通过多Agent协作与精准Prompt工程精细调优，系统响应准确率提升40\%+
  \item \textbf{项目成果：}带队获\textbf{"挑战杯"全国大学生课外学术科技作品竞赛北京市一等奖}
\end{itemize}
\datedsubsection{\textbf{创享商圈 ROBOT-XZC | 技术服务与创业实践},创始人/技术负责人}{2023 - 至今}
\begin{itemize}[parsep=0.1ex]
  \item \textbf{运营：}主导校园技术服务平台，累计完成维修及3D打印服务\textbf{2000+次}，客户满意度95\%+
  \item \textbf{培训：}策划执行机器人与3D打印公益培训\textbf{100+场}，覆盖师生\textbf{1000+人}
  \item \textbf{商业：}创新"维修盈利+硬件研发"双轮驱动，月均营收8000元
\end{itemize}

% \begin{onehalfspacing}
% \end{onehalfspacing}

% \datedsubsection{\textbf{DID-ACTE} 荷兰莱顿}{2015年}
% \role{本科毕业设计}{LIACS 交换生}
% 利用结巴分词对中国古文进行分词与词性标注，用已有领域知识训练形成 classifier 并对结果进行调优
% \begin{onehalfspacing}
% \begin{itemize}
%   \item 利用结巴分词对中国古文进行分词与词性标注
%   \item 利用已有领域知识训练形成 classifier, 并用分词结果进行测试反馈
%   \item 尝试不同规则，对 classifier 进行调优
% \end{itemize}
% \end{onehalfspacing}

\section{竞赛获奖}
\begin{multicols}{2}
\begin{itemize}[parsep=0.05ex,leftmargin=1.2pc,topsep=0.05em]
  \item "挑战杯"人工智能专项赛\textbf{北京市一等奖},2025
  \item 计算机能力挑战赛\textbf{国家级二等奖},2024
  \item 第十四届蓝桥杯\textbf{省级二等奖},2023
  \item 第十八届iCAN创新创业大赛\textbf{北京市二等奖},2024
\end{itemize}
\end{multicols}

% \section{\faHeartO\ 项目/作品摘要}
% \section{项目/作品摘要}
% \datedline{\textit{An Integrated Version of Security Monitor Vis System}, https://hijiangtao.github.io/ss-vis-component/ }{}
% \datedline{\textit{Dark-Tech}, https://github.com/hijiangtao/dark-tech/ }{}
% \datedline{\textit{融合社交网络数据挖掘的电视节目可视化分析系统}, https://hijiangtao.github.io/variety-show-hot-spot-vis/}{}
% \datedline{\textit{LeetCodeOJ Solutions}, https://github.com/hijiangtao/LeetCodeOJ}{}
% \datedline{\textit{Info-Vis}, https://github.com/ISCAS-VIS/infovis-ucas}{}


% \section{\faInfo\ 社会实践/其他}


%% Reference
%\newpage
%\bibliographystyle{IEEETran}
%\bibliography{mycite}
\end{document}
